% Folder 34, batch 04 — archivists' pages 61–80.
% Pass: Opus 5.5 (claude-opus-5-5), 2026-10-03 — first pass, unchecked against the pages by a human.
%
% Grothendieck's own pagination II.10 to II.29: the end of section 1 (abelian
% and semi-simple Galois actions on Tate modules, good reduction), section 2
% (proof of 1.1 for curves), section 3 (reduction to dimension 1 by a generic
% linear section, Lemma 3.2), section 4 (complement: H^1 mod l, almost all l)
% and the opening of section 5.
\documentclass[11pt,a4paper]{article}
\input{../preamble/grothendieck}

\folder{34}
\batch{4}
\pages{61}{80}
\foldertitle{SGA 7 : notes manuscrites (s.d.).}
\dating{[vers 1967-1973]}
\watermark{Édition de démonstration}

\begin{document}

\page{61}
\note{La page commence au milieu d'une phrase dont le début est sur la page précédente, hors de ce lot. Pagination de l'auteur : II.10 ; elle court jusqu'à II.29 en page 80.}
\uncertain{trouve} une bonne réduction sur les \struck{plus petits de} \uncertain{« plus »} de $K'$ au dessus de la base de $K$ défini par $R$). \ill{}

\textbf{Lemme 1.12.} Soit $V$ un anneau de val.\ discrète, \add{de corps résiduel $k$,} \struck{dont le corps} $\ell$ un nb premier $\neq$ car $k$, supposons \add{les notations étant celles de I~1.} que $k$ satisfasse l'hyp.\ de I~6.2. Alors pour toute représentation \add{$\ell$-adique} \uncertain{continue} abélienne de $\pi = \mathrm{Gal}(\overline{K}/K)$, la restriction à $I(\ell)$ est \uncertain{essentiellement} triviale.

Se déduit de I~6.2 comme \struck{\ill{}} on vient de déduire \struck{1.9} 1.4 de 1.1 via 1.2.
\marginal{Pour le cas des \ill{} de car.~$\ell$, cf feuille à la fin des présentes notes p.~13}\note{note encerclée, reliée par un trait à « Remarque 1.13 » ; « p. 13 » est la pagination de l'auteur.}

\textbf{Remarque 1.13.} Le seul exemple connu où on sache vérifier l'hypothèse \struck{\ill{}} de l'action \struck{abélienne} \add{abélienne et} (semi-simple) d'un groupe de Galois sur $T_\ell(A_{\overline{K}})$ est celui où $A$ est une variété « de type CM » au sens de Shimura--Taniyama [~]. \add{Nous voyons donc que} \struck{Dans ce cas} (le corps de définition de $A$ étant quelconque) $A$ devient radiciellement isogène, sur une extension \struck{finie} séparable finie convenable, à un schéma abélien définissable sur un corps fini (\uncertain{cas} car.~$p>0$), resp.\ sur l'anneau des entiers d'un corps de nombres (cas de la car.~$p=0$). Ce dernier cas \add{tout au moins} étant bien connu [~].

Il semble plausible qu'en caractéristique

\page{62}
nulle, \struck{\ill{}} un schéma abélien sur un corps de type fini $K$ (\uncertain{qu'on peut} $=$ supposer un corps de nombres) tel que l'action de $\mathrm{Gal}(\overline{K}/K)$ sur $T_\ell(A(\overline{K}))$ soit \struck{ess.} abélienne \add{(ou seulement ess.\ résoluble)}, soit de type CM.

1.14. On notera cependant que \struck{les} \add{certains des} arguments donnés dans le présent numéro dépassent le cadre des schémas abéliens, et s'appliquent \struck{\ill{}} \add{\ill{}} à des « motifs » de \struck{degré} \add{poids} quelconque (une fois la théorie des motifs assise sur des bases solides). Notons seulement qu'une démonstration de réduction à une situation arithmétique, ess.\ identique à la démonstration de 1.5, permet de prouver ceci :

\textbf{Théorème 1.15.} Soient $k$ un corps alg.\ clos, $Y$ un schéma \struck{\ill{} de t.f.\ sur $k$, sur $k$, \ill{}} loc.\ de t.f.\ \add{connexe} et normal, $\xi$ un pt géométrique \struck{\ill{}} de $Y$, $f \colon X \to Y$ un morphisme propre \struck{\ill{}} \struck{\ill{}} \add{\ill{}}, $\ell$ un nb premier $\neq$ car.~$k$, \add{$i$ un entier,} supposons que \struck{la représentation} $R^i f_*(\mathbb{Z}_\ell)$ soit un faisceau $\ell$-adique \uncertain{constant}\marginal{(\ill{} toujours vrai à condition de \uncertain{rapetisser} $Y$)}. Considérons la représentation de $\pi_1(X,\xi)$\note{il écrit $\pi_1(X,\xi)$ ici et deux lignes plus bas ; la représentation est celle du groupe fondamental de la base $Y$.} dans $E = H^i(X_\xi, \mathbb{Q}_\ell)$, \add{et le groupe algébrique $G \subset \mathrm{Aut}_{\mathbb{Q}_\ell}(E)$ engendré.} Alors le radical $R$ de $G^0$ est unipotent. \struck{\ill{} cette représentation est \ill{} ess.} En particulier, si \struck{\ill{} elle est \ill{}} la représentation envisagée de $\pi_1(X,\xi)$ est ess.\ résoluble, elle est ess.\ \struck{\ill{}} \struck{et de \ill{}} \add{unipotente ; donc si on la suppose de plus semi-simple (*), elle} est essentiellement triviale.

1.16. Dans l'énoncé 1.15, l'hypothèse de \uncertain{propreté} ne nous a servi qu'à assurer que \struck{la formation de} $R^i f_*(\mathbb{Z}_\ell)$ est un faisceau $\ell$-adique \emph{constructible}, au voisinage du point générique $y$ de $Y$.

(*) J'ai \struck{\ill{}} \ill{} dans 1.5 la conjecture que \add{si $f$ est lisse,} l'hypothèse de ½simplicité est automatiquement vérifiée.\note{« ½simplicité » : semi-simplicité, dans son abréviation habituelle.}

\page{63}
\struck{constante tordue sur $X_y$, dont la formation commute à \ill{}, \ill{} dans la formation commute à \ill{}}
\ill{} que l'énoncé analogue est vrai après réduction de la situation à un corps de base de type fini. Cette hypothèse est \add{donc} \uncertain{vérifiée} si on suppose qu'on dispose de la résolution des singularités pour les schémas de type fini \add{de dim $\leq \dim X_y$} sur une extension algébrique de $K$ ($y$ étant le pt générique de $Y$), \struck{D} compte tenu I~8.8, 8.9]. Il en est ainsi en particulier si car $k = 0$, ou si $\dim X_y \leq 2$.

On peut montrer que la conclusion est également valable si $i = 1$, sans restriction \struck{\ill{}} de propreté sur $f$, \struck{\ill{} cf} en utilisant les réductions faites dans les nos suivants, comparer n°~5.

\note{La remarque qui suit est barrée de deux longs traits obliques ; elle se lit néanmoins.}
\struck{\textbf{Remarque 1.17.} Il est plausible que l'on ait} \struck{Gardons les notations de 1.1, \struck{\ill{}} et soit $H$ un groupe \add{profini} quotient de $I$, invariant par l'action de $\pi$ par automorphismes \struck{\ill{}}, et tel que $H$ soit un \struck{\ill{}} pro-$\ell$-groupe \emph{résoluble}. On trouve donc une extension}
\[
  \text{\struck{$1 \to H \to \pi' \to \pi_0 \to 1$}} .
\]
\struck{Ceci posé, il est plausible que l'on ait \ill{} $[\pi_0,H]$ \ill{} de $\pi'$. Cela semblerait devoir résulter de l'hyp.\ de Riemann--Weil en dimension quelconque, comme 1.1 résulte de l'hyp.\ de R.W. en dim~1. S'il en était ainsi, on}

\page{64}
\section*{Fin de la démonstration de 1.9}

\emph{Bonne réduction \struck{\ill{} de car.} (action ess.\ nilpotente).}

a) L'action est semi-simple et abélienne.

On utilise les faits suivants :

\textbf{Lemme.} Si $A$ est un schéma abélien sur un corps fini $k$, l'action des \add{relèvements des} Frobenius sur $T_\ell(A(\overline{k}))$ est semi-simple.

En effet, \struck{\ill{}} cette action est \struck{\ill{}} \add{\uncertain{induite par}} celle du Frobenius géométrique \struck{\ill{}}, endomorphisme de $A_{\overline{k}}$\note{l'indice de $A$ est surchargé.}. \struck{Donc} Or cet \uncertain{endomorphisme} appartient au \emph{centre} de l'anneau $\mathrm{End}(A) \otimes_{\mathbb{Z}} \mathbb{Q}$ des endomorphismes \add{tensorisé par $\mathbb{Q}$}, qui, par Weil [2], est une algèbre \add{finie} semi-simple sur $\mathbb{Q}$, dont le centre $Z$ \struck{\ill{}} est une algèbre finie \struck{\ill{}} \add{étale} de $\mathbb{Q}$. Donc $Z_\ell = Z \otimes_{\mathbb{Q}} \mathbb{Q}_\ell$ \struck{\ill{}} est une algèbre finie étale sur $\mathbb{Q}_\ell$, \struck{Donc $Z_\ell$ opère} \ill{} $Z$ opère de façon semi-simple sur $T_\ell(A) \otimes_{\mathbb{Z}_\ell} \mathbb{Q}_\ell$.

\textbf{\struck{Lemme} \add{Corollaire 1.}} Soit alors $S$ un schéma \add{loc.}\ de t.f.\ sur $\mathrm{Spec}\,\mathbb{Z}$, $\xi$ un pt géom.\ de $S$, $A$ un schéma abélien sur $S$, \add{$\ell$ un nb premier distinct des car.\ résiduelles de $S$,} d'où une opération de $\pi_1(S,\xi)$ sur \add{$E =$} $T_\ell(A(\xi)) \otimes_{\mathbb{Z}_\ell} \mathbb{Q}_\ell$. Pour tout point fermé $s \in S$, les éléments de Frobenius $f_s \in \pi_1(S,\xi)$ [déterminés à \struck{\ill{}} conj.\ près comme l'image du gén.\ de Frob.\ de $\pi_1(s)$ par $\pi_1(s) \to \pi_1(S)$] opèrent \struck{\ill{}} de façon semi-simple.

\textbf{Corollaire 2.} Si la représentation de $\pi_1(S,\xi)$ sur $E$ est ess.\ nilpotente, elle est \struck{abélienne et} semi-simple et ess.\ \uncertain{abélienne}.

\page{65}
En effet, le groupe \add{$G \subset \mathrm{Aut}_{\mathbb{Q}_\ell}(E_\ell)$} algébrique (engendré par cette représentation) est tel que $G^0$ est nilpotent, \struck{D'autre} \struck{\ill{}} donc : produit $T \times U$, $T$ un tore, $U$ unipotent. \add{À prouver $U = 1$.} Soit $\pi' \subset \pi$ l'image inverse de $G^0$ ; il correspond à un revêtement \struck{fini} étale connexe $S'$ de $S$, et quitte à passer à celui-ci, on peut supposer $G = G^0$. On utilise le th.\ de \emph{Čebotarev} \struck{\ill{}} que les \emph{éléments de Frobenius sont denses dans $\pi_1(S,\xi)$}. On en conclut que leurs images $f_s$ dans $G$ engendrent $G$ comme groupe algébrique. Or $T$ est \uncertain{V} \ill{} des éléments semi-simples de $G$, donc $G = T$, i.e.\ $U = (1)$, cqfd.

b) \emph{Fin de la démonstration.} Je \uncertain{dis} d'abord que \struck{\ill{}} pour tout $u \in \mathrm{Max}(R)$ tel que $R/u$ \uncertain{soit} de car.~$\ell$, $\exists$ extension finie $K'$ de $K$ telle que $A_{K'}$ ait bonne réduction en tous les idéaux maximaux de $V'$ (où $V'$ est le normalisé de $V = R_u$ dans $K'$). Pour ceci, \struck{l'action} on note que \struck{\ill{}} en vertu de cor.~2, l'action de

\note{La fin de la page est barrée d'un trait horizontal et encadrée de deux longs traits obliques ; elle se lit ainsi.}
\struck{utilisant un résultat de \ill{} cap.\ ultérieur, \ill{} sans \ill{} utilisant le fait que \struck{l'action} \ill{} \ill{} extension \ill{} finie $K'/K$ telle que \struck{l'extension} le modèle de Néron de $A_{K'}$ sur $R'$ ait « réduction ½stable » \uncertain{tout} \struck{au point} \ill{} de $V'$, \ill{} $V'$. \struck{Alors \ill{}} Plus…}
$\mathrm{Gal}(\overline{K}/K)$ sur $T_\ell(A(\overline{K})) \otimes_{\mathbb{Z}_\ell} \mathbb{Q}_\ell$ est ½simple abélienne, donc \uncertain{est} \ill{} \uncertain{a fortiori sur celles des} groupes d'inertie $I_u$.

\page{66}
\struck{choisissant $K'$ assez grand, et utilisant le cor.~2 précédent, \ill{} que \uncertain{l'action} de $\mathrm{Gal}(\overline{K'}/K')$ \ill{} abélienne. On garde les notations, on peut supposer \uncertain{alors} $K' = K$. Or l'action \ill{}}\note{Les quatre premières lignes sont encadrées, barrées et biffées en oblique.}

Je \uncertain{me suffis} d'établir le résultat suivant (qui a été signalé par Serre) :

\textbf{Lemme.} Soient $V$ anneau de val.\ discrète, corps de fractions $K$, $\ell$ un nb premier distinct de la car.\ de $K$ (\uncertain{mais pas nécessairement} de la car.\ du corps résiduel), $A$ schéma abélien sur $K$, \struck{\ill{}} supposons que l'action du groupe d'inertie \struck{\ill{}} $I \subset \mathrm{Gal}(\overline{K}/K)$ sur $E = T_\ell(A(\overline{K})) \otimes_{\mathbb{Z}_\ell} \mathbb{Q}_\ell$ soit \struck{\ill{}} \add{\ill{}} semi-simple, \struck{\ill{} abélienne, \ill{}} alors \struck{$A$ a bonne} \ill{} (\uncertain{ceci} étant \uncertain{sous-entendu}) $A$ a bonne réduction « potentielle ».

\marginal{Pour la démonstration, renvoyer à exposé ultérieur}
Le cas où la car.\ résiduelle $\neq \ell$ est évident, car l'action de $I$ étant ess.\ unipotente (\uncertain{J}), \struck{et} est semi-simple, \uncertain{si elle} doit être ess.\ triviale. \struck{\ill{}} Le cas où la car.\ résiduelle est $\ell$ (dans $K$ de car.\ zéro) demande un traitement \uncertain{\ill{}} \struck{\ill{}} \ill{} on peut supposer $V$ \struck{\ill{}}, \struck{et complet, \ill{}} \add{à corps résiduel sép.\ clos.} Considérons le modèle de Néron $G$ de $A$ sur $V$ [~]. \struck{Alors} Le th.\ de réduction ½stable, qui se \uncertain{prouve} ultérieurement, indique que (quitte à faire une extension finie sur $K$) on peut supposer que $G^0_k$ est ext.\ d'une V.A.\ par un tore \add{$T$ de dim.~$r$}. On sait alors que $T_\ell(A(\overline{K}))$ est extension d'un $\mathbb{Z}_\ell$-

\page{67}
module $M \simeq \mathbb{Z}_\ell^r$ \struck{\ill{}} sur lequel $I$ opère trivialement, \add{canoniquement : $T_\ell(G(\overline{k}))$,} par un \struck{groupe} \add{$N$} abélien de rang $2g-r$, \uncertain{bonnement}. \uncertain{Si} l'action \struck{de} de $I$ est semi-simple, il s'ensuit que l'on a un sous-module $M \simeq \mathbb{Z}_\ell^r \subset T_\ell(A(\overline{K}))$ \emph{sur lequel} $I$ \emph{opère trivialement} \struck{\ill{}}, \add{avec $M \cap N = 0$. On en déduit} (un sous-groupe de la forme $(\mathbb{Q}_\ell/\mathbb{Z}_\ell)^r$ du groupe $A(K)$ ; par définition du modèle de Néron, il serait contenu dans $G(k)$, i.e.\ $M \subset N$, donc $M = 0$, donc $r = 0$, ce qui signifie que la réduction est stable, cqfd.

\textbf{Remarque de Serre.} Si on savait montrer que (dans le cas de la car.\ nulle) une action résoluble (\uncertain{indique} \add{(VA sur un corps de type fini)}) action ess.\ abélienne, \ill{} on pourrait en conclure qu'une \add{telle} VA est « de type CM » après extension finie du corps de base.

On est ramené, par ce qui précède, au cas \struck{d'une VA} \add{d'un schéma abélien} sur le \uncertain{normalisé} \add{$R$} d'un corps de nbs $K$, avec action abélienne de $\mathrm{Gal}(\overline{K}/K)$ sur $E_\ell = T_\ell(A(\overline{K})) \otimes \mathbb{Q}_\ell$, i.e.\ de $\pi_1(S,\xi)$ ($S = \mathrm{Spec}\, R$). L'action de $\pi_1(S,\xi)$ est loc.\ algébrique (conditions locales \add{en tout $u \in \mathrm{Max}(R)$}) sur l'action des $\pi_1$ des corps locaux, i.e.\ des groupes de \struck{\ill{}} décomposition) et « \add{loc.}\ rationnelle », i.e.\ \add{p.\ ex.\ les pol.\ car.\ des \uncertain{Frobenius sont à coeff} $\in \mathbb{Q}$}. Or on peut montrer que si on a deux telles représentations, dans un $E_\ell$ \emph{et}

\page{68}
dans un $E_{\ell'}$, l'une étant semi-simple abélienne, \struck{alors} \struck{l'autre est} ayant les mêmes pol.\ car.\ de Frob., pour presque tous $u \in \mathrm{Spec}\, R$, alors l'autre représentation est \emph{résoluble},\note{Un passage de quatre lignes est ici encadré, barré horizontalement et biffé de traits obliques ; on y devine $E_\ell$, $E_{\ell'}$, $T_\ell$, $T_{\ell'}$, « ess.\ résoluble », « la représentation ½simple » et « représentation sur $\mathbb{Q}$ ».} \struck{\ill{} aussi \ill{} abélienne. Alors \ill{} \ill{}} et pour \ill{} $\mathrm{Spec}\, R$, la repr.\ \emph{½simple associée} à $E_{\ell'}$ \struck{\ill{}} est la \uncertain{même} que celle qui est définie par $E_\ell$. Choisissons alors $\ell$ tel que
\note{La page s'arrête sur cette phrase inachevée ; le reste de la feuille est blanc.}

\page{69}
\note{Le haut de la page, barré de traits obliques, se lit ainsi.}
\struck{\ill{} pourrait généraliser substantiellement 1.3, 1.5, \ill{} 1.9 (sous réserve de la ½simplicité), 1.13, en y remplaçant partout « ess.\ abélienne » par « ess.\ résoluble ».}

\section*{2. Démonstration de 1.1 : cas où $X$ est de dim~1}

\struck{Alors $X$ est une courbe}

2.1. Notons d'abord (sans hyp.\ sur la dim.\ de $X$\struck{, \ill{} \ill{}}) qu'on peut supposer $X$ \emph{affine}. En effet, si $U$ est un ouvert affine \struck{\ill{}} de $X$ contenant l'image de $\xi$, on a \add{\ill{} de suites exactes} \struck{\ill{} que $\pi_1(\ill{})$}

\begin{tikzcd}[column sep=small]
1 \arrow[r] & \pi_1(\overline{U},\xi) \arrow[r] \arrow[d, "\alpha"] & \pi_1(U,\xi) \arrow[r] \arrow[d, "\beta"] & \pi_1(S,\xi) \arrow[r] \arrow[d, "\wr"] & 1 \\
1 \arrow[r] & \pi_1(\overline{X},\xi) \arrow[r] & \pi_1(X,\xi) \arrow[r] & \pi_1(S,\xi) \arrow[r] & 1
\end{tikzcd}

\note{L'étiquette de la flèche verticale de droite est un trait doublé, lu comme un isomorphisme.}
Comme \struck{\ill{}} $X$ \struck{est} \add{\uncertain{normal},} \struck{\ill{}} \struck{\ill{} $\overline{X}$ cf.\ EGA~IV 6.15.6}, il s'ensuit que $\beta$ est surjectif (SGA~1~V~…), donc l'est aussi $\alpha$. \struck{$\pi_1(X)$ \ill{} $\alpha$ et $\beta$ sont surjectifs.} On en conclut que si les \struck{\ill{}} conclusions de 1.1 \uncertain{sont} vraies pour $(U,\xi)$, elles sont vraies pour $(X,\xi)$.

\marginal{2.2. On voit que la condition \ill{}}\note{Note marginale de neuf lignes écrite de biais dans la marge gauche, barrée de deux traits ; on y devine $K'$, $X(\ill{})$, $\pi_1(X,\xi)$, $\pi_1(S,\xi)$, le reste est illisible.}

2.3. \struck{Sans} \struck{Alors} Supposons donc \struck{$X$} affine, donc\note{La suite, jusqu'au bas de la page, est biffée de traits obliques ; elle se lit néanmoins.}
\struck{quasi-projectif, et supposons de plus $\dim X = 1$. Alors $X$ est un \struck{\ill{}} ouvert d'une courbe normale projective $\hat{X}$ (EGA~II~7.…). Quitte à faire une extension \struck{radicielle} finie de $K$, on est ramené au cas où $\hat{X}$ est \uncertain{une} \emph{lisse} sur $K$, et où les points de $\hat{X} - X$ sont \uncertain{tous} rationnels sur $K$. Soient $x_i$ ($1 \leq i \leq n$) ces pts.}

\page{70}
\note{Cette page reprend, sous une autre forme, les nos 2.2 et 2.3 de la page précédente.}
\struck{On a alors la suite exacte bien connue :}

2.2. On voit de même que la conclusion ne change pas \struck{$\pi_1(X$} si on remplace $k$ par une extension finie $k'$, $X$ par $X' = X \otimes_k k'$, et $\xi$ par un pt géom.\ $\xi'$ au-dessus de $\xi$. En effet, cela revient à remplacer $\pi_1(X,\xi)$ par un s.-groupe ouvert, image inverse du s.-groupe ouvert de $\pi_1(S,\xi)$ défini par l'extension $k'/k$.

2.3. Supposons maintenant $X$ quasi-projectif, comme il est loisible par 1.1. Il existe alors un schéma \struck{\ill{}} projectif $X'$, sur $k$, tel que $X$ soit \struck{\ill{}} \add{dense} un ouvert de $X'$. On peut de plus supposer $X'$ normal (quitte à le remplacer par $(X'_{\mathrm{red}})$ \add{normalisé de}). Utilisant \struck{\ill{}} EGA~IV 17.15.14.2, on voit qu'il \struck{existe} existe une extension finie radicielle $k_1$ de $k$, telle que $X_1$ \add{normalisé de} $X'_1 = X' \otimes_k k_1$ soit géom.\ normal. \struck{Il \ill{}} \ill{} \ill{} \ill{} des \ill{} \ill{} $X \otimes_k k_1$, qui s'identifie à \struck{\ill{}} un ouvert de $\widetilde{X}'_1$. \struck{\ill{}} Comme $X_1$ est géom.\ unibranche (EGA~IV 6.15.6), $\overline{X}_1 \to X_1$ est radiciel surjectif, donc induit un \uncertain{isom.}\ sur \add{\ill{} $\overline{X}_1 \to \overline{X'_1}$} les $\pi_1$ (SGA~1~VIII …). Donc, \uncertain{conjuguant} cette remarque avec 2.2, on est ramené au cas où $X'$ est géom.\ normale.

\page{71}
2.4. On sait déjà que $V = \pi_1(\overline{X})^{\mathrm{ab}}(\ell)$ est un module de t.f.\ sur $\mathbb{Z}_\ell$, \struck{ce qui} revient au même que $H^1(\overline{X}, \mathbb{Z}/\ell\mathbb{Z})$ est un groupe fini. Ceci est vrai pour $X$ schéma de t.f.\ sur un corps \add{sép.}\ \struck{alg.}\ clos, et \add{la} démonstration par réduction au cas des courbes, \uncertain{on} la prouvera dans no~3. Lorsqu'on \add{sait} \struck{\ill{}} que $\pi_1(\overline{X})^{\mathrm{ab}}(\ell)$ est de t.f.\ sur $\mathbb{Z}_\ell$, dans la relation : \struck{\ill{}}\note{dans la relation énoncée en 1.1, vraisemblablement.} prouvons, \struck{\ill{} \ill{}} \struck{$\pi_1(\overline{X})(\ell)_{\pi_0}$ fini i.e.\ \ill{}} \struck{(\ill{})} $V_{\pi_0}$ fini i.e.\ $(V \otimes_{\mathbb{Z}_\ell} \mathbb{Q}_\ell)_{\pi_0} = 0$, et par dualité elle \uncertain{équivaut} à la relation
\[
  (\check{V})^{\pi_0} = 0
\]
où $\check{V} = \mathrm{Hom}_{\mathbb{Z}_\ell}(V, \mathbb{Z}_\ell)$. Or on a un isom.\ canonique $\check{V} \simeq \mathrm{Hom}(\pi_1(\overline{X},\xi), \mathbb{Z}_\ell) \simeq H^1(\overline{X}, \mathbb{Z}_\ell)$, évidemment compatible avec l'action de $\pi_0$. Donc la conclusion de 1.1 s'écrit de façon équivalente, \uncertain{en} termes cohomologiques
\begin{equation}
  H^1(\overline{X}, \mathbb{Z}_\ell)^{\pi_0} = 0 . \tag{2.4.1}
\end{equation}

\marginal{faire une section 2.4…\ \uncertain{spéciale} \ill{} \struck{\ill{}} être \uncertain{vers}}
\emph{Faire} d'abord une réduction au cas où $\pi = \pi_1(X,\xi)$, en introduisant \struck{\ill{}} le rev.\ étale fini $X'$ de $X$ qui correspond à $\pi \subset \pi_1(X,\xi)$.\note{Ces deux dernières lignes sont encadrées, reliées à la note marginale.}

\page{72}
2.5. Dans ce qui suit, nous supposons $\dim X = 1$, \add{$\dim X' = 1$}. Quitte à faire encore une extension finie sur $k$, on peut supposer les points $x_i$ ($1 \leq i \leq n$) de $X' - X$ rationnels sur $k$. On a alors une suite exacte bien connue
\[
  0 \to H^1(\overline{X'}, \mathbb{Z}_\ell) \to H^1(\overline{X}, \mathbb{Z}_\ell) \to \prod_{1 \leq i \leq n} H^2_{x_i}(\overline{X}, \mathbb{Z}_\ell)
\]
\[
  \to H^2(\overline{X'}, \mathbb{Z}_\ell) \to 0 \tag{2.5.1}
\]
\note{La flèche finale vers $0$ est tracée verticalement sous $H^2(\overline{X'},\mathbb{Z}_\ell)$.}
avec les flèches \uncertain{sont} \struck{\ill{}} (évidemment) compatibles avec l'action de \struck{\ill{}} $\pi_0 = \mathrm{Gal}(\overline{k}/k)$. De plus, on a des isomorphismes canoniques
\[
  H^2_{x_i}(\overline{X}, \mathbb{Z}_\ell) \simeq \mathbb{Z}_\ell(-1) \tag{2.5.2}
\]
\note{À la suite de cette formule, une expression biffée : $\varprojlim_\nu {}_{\ell^\nu}\overline{k}^*$.}
où l'on a
\[
  \mathbb{Z}_\ell(-1) = \mathbb{Z}_\ell(1)^{\otimes(-1)}, \qquad \mathbb{Z}_\ell(1) = \varprojlim_\nu \left( {}_{\ell^\nu}\overline{k}^* \right) . \tag{2.5.3}
\]
(De plus, on a un isom.\ canonique $H^2(\overline{X'}, \mathbb{Z}_\ell) \simeq \mathbb{Z}_\ell(-1)$, et l'homomorphisme $\prod \to H^2(\overline{X'}, \mathbb{Z}_\ell)$ s'identifie à l'homomorphisme somme.) L'isomorphisme (2.5.2) est encore compatible avec l'action de $\pi_0$, opérant sur le second membre de façon explicitée par la formule (2.5.3). \struck{\ill{}}

Il s'ensuit que en tant que $\pi_0$-module $\ell$-adique, $H^1(\overline{X}, \mathbb{Z}_\ell)$ est une extension de $\mathbb{Z}_\ell(1)^{n-1}$ par $H^1(\overline{X'}, \mathbb{Z}_\ell)$. Comme l'opération de $\pi_0$ sur $\mathbb{Z}_\ell(1)$ est non triviale, il s'ensuit \add{déjà} qu'on a
\note{Le signe de la torsion : il écrit $\mathbb{Z}_\ell(-1)$ en (2.5.2) et $\mathbb{Z}_\ell(1)^{n-1}$ ici, comme sur la page.}

\page{73}
\[
  H^1(\overline{X}, \mathbb{Z}_\ell)^{\pi_0} = H^1(\overline{X'}, \mathbb{Z}_\ell)^{\pi_0} . \tag{2.5.4}
\]
\struck{Je vais \add{\ill{}} \ill{}} \add{\uncertain{Reste} à voir} que le dernier membre est nul :
\[
  H^1(\overline{X'}, \mathbb{Z}_\ell)^{\pi_0} = 0 . \tag{2.5.5}
\]
C'est une conséquence connue de l'hyp.\ de Riemann--Weil, prouvée en dim.~1. Explicitons l'argument. \struck{Soit $A$ \ill{}} Posons $\overline{X'} = C_k$.\note{lecture douteuse de ce qui suit « Posons ».} \struck{\ill{} On sait que $k$} \ill{} \ill{} \struck{\ill{}} \add{corps} des fonctions est une \struck{\ill{}} \add{\uncertain{extension}} \ill{} \ill{} de type fini sur $\mathbb{Z}$, \struck{\ill{}} telle que \ill{} $\boxtimes$ \uncertain{$C$} \uncertain{provienne} d'un schéma $C_S$ sur $S = \mathrm{Spec}\,R$, propre et lisse sur $S$ (\add{on \uncertain{peut supposer} $S$ normal}). \struck{\ill{}} La \add{relation} \struck{\ill{}} (2.5.5) s'écrit alors
\[
  H^0(S, R^1 f_*(\mathbb{Z}_\ell)) = 0
\]
\add{\ill{}} où $f \colon C \to S$. Sous cette forme, on voit \struck{\ill{}} que cette relation \struck{\ill{}} ne change pas si on remplace $S$ par \struck{\ill{}} un ouvert \ill{} \struck{\ill{}} \add{non vide} de $S$, donc \ill{} \ill{} \ill{} \ill{} \uncertain{ponctuel} $\xi$ de $S$, par \struck{\ill{}} \add{\ill{}} l'opération de $\pi_1(S,\xi)$ sur $H^1(C_{\overline{\xi}}, \mathbb{Z}_\ell)$, \struck{\ill{}} \add{pour $\xi$ choisi, on \uncertain{pourra} supposer} \ill{} de $S$, donc \ill{} \ill{} \uncertain{réduire} \struck{Alors} à prouver, \ill{} \ill{} se \uncertain{réduit} à la relation
\[
  H^1(C_{\overline{s}}, \mathbb{Z}_\ell)^{\pi_1(s,\overline{s})} = 0 \tag{2.5.6}
\]
pour \struck{\ill{}} \ill{} fini.\note{La rédaction de ce paragraphe est très surchargée ; seules les formules et l'articulation se lisent avec sûreté.}

Ceci, ramène au cas d'un corps de base fini. Mais dans ce cas on sait \add{\uncertain{par} la \uncertain{généralité}, il \uncertain{suffira} donc que} $\pi_0 = \hat{\mathbb{Z}}$, \struck{\ill{} engendré \ill{} de Frobenius}
\struck{\ill{}, \ill{} Frob de $\hat{\mathbb{Z}}$ opérant sur $H^1(C_{\overline{s}}, \mathbb{Z}_\ell)$ [2, …] \ill{} \ill{} valeurs absolues $q^{1/2}$ ($q = \mathrm{card}\,k$). Donc elles sont $\neq 1$ \ill{} (2.5.6)} \ill{} ce qui est trivial.\note{Ce passage, encadré, est biffé de traits obliques.}

\marginal{\textbf{Remarques 2.6.} Nous avons utilisé en \uncertain{substance} 2 propriétés \uncertain{du} corps $K$ : a) \struck{\ill{}} \ill{} $\in \overline{\pi}$ \ill{} les racines $\ell^\nu$-ièmes de l'unité, \ill{} est infini. b) Pour \ill{} schéma abélien $A$ sur une extension finie $K'$ de $K$, $T_\ell(A)_\pi$ est fini, où $\pi = \mathrm{Gal}(\overline{K'}/K')$ \ill{} \uncertain{ce qui revient à} \uncertain{savoir} que $T_\ell(A)^\pi$ \ill{}. Pour b) il ne \uncertain{suffit} \ill{} \ill{} (exemple : $K = \mathbb{R}$\ill{}).}
\marginal{\struck{$\mathrm{Ker}(1-f)$} dans $H^1(C_{\overline{s}}, \mathbb{Z}_\ell) \simeq T_\ell(A)^\vee$ est nul, i.e.\ $\mathrm{Coker}(1-f)$ dans $T_\ell(A)$ fini, i.e.\ $\mathrm{Ker}(1-f)$ dans $T_\ell(A)$ fini, i.e.\ \ill{} $A(k)$ fini, ce qui est trivial.}
\note{Deux notes écrites en long dans la marge gauche ; dans un cartouche, en haut, des renvois biffés illisibles (« EGA » s'y devine).}

\page{74}
\section*{3. Réduction au cas de la dimension 1}

3.1. \struck{Supp.} Soit d'abord $k$ un \struck{\ill{}} corps quelconque, $X'$ un schéma \add{projectif \uncertain{intègre} \uncertain{normal}, géom.\ connexe} \struck{géométriquement unibranche} sur $k$, $X$ un ouvert dense de $X'$\add{, $\xi$ un pt géom.\ de $X'$, $\overline{k}$ la clôture} séparable de $k$ dans $k(\xi)$, $\overline{X'}$, $\overline{X}$ déduits de $X'$, $X$ par $k \to \overline{k}$. \struck{On a \ill{} \ill{}} On veut étudier $\pi_1(\overline{X},\xi)$, et \struck{\ill{}} \ill{} nous nous intéressons en fait au groupe fondamental d'une courbe lisse sur un corps convenable. Comme on sait \add{en \ill{}}, $\pi_1(\overline{X},\xi)$ \uncertain{est} \uncertain{muni} de \uncertain{façon} naturelle d'un \uncertain{groupe} d'opérateurs, \uncertain{à} savoir $\pi_1(S,\xi)$ ($S = \mathrm{Spec}\,k$), qui \uncertain{permet} de \uncertain{considérer} aussi la suite exacte
\[
  1 \to \pi_1(\overline{X},\xi) \to \pi_1(X,\xi) \to \pi_1(S,\xi) \to 1 . \tag{3.1.1}
\]
\struck{Soit $T$ le schéma \ill{}} Nous supposons $X'$ immergé dans un $\mathbb{P}^r_k$\struck{\ill{}}. Soit $n = \dim X'$, et soit $T$ la grassmannienne des \uncertain{sous-espaces} \uncertain{linéaires} (EGA~I, 2\ill{} \uncertain{édition}, …) \add{projectifs} de codimension $n-1$ dans $P$, $t$ le point générique de $T$, $K = k(t)$, qui est une extension \add{\uncertain{de t.f.}}\ \struck{pure} de $k$. Soit $L_t$ le \uncertain{sous-espace} \add{« générique »} de codimension $n-1$ de $P_K$ qui correspond à $t \in T_K$, et considérons
\[
  C' = L \cap X'_K .
\]
Il est connu que $\dim C' = 1$, que $C'$ est géom.\ \add{\uncertain{connexe}} normal et géom.\ connexe. Soit $C$ l'image inverse de $X_K$ dans $C'$.

\begin{tikzcd}
C \arrow[r, hook] \arrow[d] & C' = L_t \cap X'_K \arrow[d, no head] \\
X_K \arrow[r, hook] & X'_K
\end{tikzcd}

\note{Petit diagramme écrit dans la marge gauche, en face de la définition de $C'$ ; la flèche verticale de droite est un simple trait.}

\page{75}
Considérons \struck{l'image \ill{}} le morphisme canonique
\[
  C \to X
\]
composé de l'inclusion $C \to X_K$ et $X_K \to X$. Nous supposons que \struck{\ill{}} l'image de $\xi$ dans $X$ est contenue dans celle de $C$. Ceci signifie \struck{\ill{} \ill{}} \uncertain{seulement} l'exclusion du cas (\uncertain{exclu} \ill{} $\dim \overline{\{x\}} \geq 1$) ; \struck{\ill{}}

3.1.1. \uncertain{Ceci} équivaut, à \uncertain{\ill{}} près, dans le cas d'\uncertain{hypothèse} \uncertain{anodine}. Soit $\xi$, cela est une \uncertain{\ill{}} $\xi_1$ un point géométrique de $C$ au-dessus de $\xi$,
\[
  \begin{array}{ccc}
    \xi_1 & \longrightarrow & \xi \\
    \downarrow & & \downarrow \\
    C & \longrightarrow & X \\
    \downarrow & & \downarrow \\
    \mathrm{Spec}\,K & \longrightarrow & S
  \end{array} \tag{3.1.2}
\]
et soit $\overline{K}$ la clôture sép.\ de $K = k(t)$ dans $k(\xi_1)$.
\note{Un astérisque est placé à droite de cette phrase, sans renvoi visible.}

\note{Le passage qui suit est encadré et biffé de deux traits obliques.}
\struck{On a un diagramme commutatif d'inclusions de corps}
\[
  \begin{array}{ccc}
    k & \longrightarrow & K = k(t) \\
    \downarrow & & \downarrow \\
    \overline{k} & \longrightarrow & \overline{K}
  \end{array}
\]
\note{Ce diagramme est biffé avec le passage qui l'entoure.}
\struck{Considérons aussi le corps $K\overline{k}$ \uncertain{composé} dans $\overline{K}$, qui n'est autre que le corps $k(\overline{t})$, où $\overline{t}$ est le point générique de $T_{\overline{k}}$.}

Le diagramme (3.1.2) donne un diagramme commutatif, à lignes exactes :

\page{76}
\begin{tikzcd}[column sep=small]
1 \arrow[r] & \pi_1(\overline{C},\xi_1) \arrow[r] \arrow[d, "\alpha"] & \pi_1(C,\xi_1) \arrow[r] \arrow[d, "\beta"] & \pi_1(\mathrm{Spec}\,K,\xi_1) \arrow[r] \arrow[d, "\gamma"] & 1 \\
1 \arrow[r] & \pi_1(\overline{X},\xi) \arrow[r] & \pi_1(X,\xi) \arrow[r] & \pi_1(\mathrm{Spec}\,k,\xi) \arrow[r] & 1
\end{tikzcd}

\note{Diagramme numéroté (3.1.3). Au-dessus de $\pi_1(\mathrm{Spec}\,K,\xi_1)$, écrit sur un terme raturé, il note « $\mathrm{Gal}(\overline{K}/K)$ » relié par un signe d'égalité.}
où on a posé $\overline{C} = C \otimes_K \overline{K}$. \struck{Ceci posé, je}

Ceci posé, on a

\textbf{Lemme 3.2.} Avec les \add{hyp.\ et} notations précédentes, dans le diagramme (3.1.3), les \add{trois} \add{homomorphismes} \struck{\ill{}} \add{verticaux} \struck{$\alpha$ \ill{} \ill{}} sont \emph{surjectifs}.

\struck{\ill{} Pour prouver le lemme}, Il suffit de prouver la surjectivité de $\gamma$ et $\alpha$. Celle de $\gamma$ provient du fait que $K = k(t)$ est une extension \struck{transcendante} pure de $k$, donc \struck{\ill{}} $k$ est algébriquement fermé dans $K$. Notons maintenant que la flèche $\overline{C} \to \overline{X}$ donnant naissance à $\alpha$ est la même que celle qu'on obtiendrait en \struck{remplaçant \ill{} \ill{}} faisant le changement de corps de base $k \to \overline{k}$. Donc on peut supposer $k$ séparablement clos, et \struck{\ill{}} on \struck{\ill{} \ill{} à \ill{}} \add{doit} examiner le morphisme sur les $\pi_1$ induit par
\[
  C \to X .
\]
Dire que ce morphisme sur les $\pi_1$ est surjectif signifie que pour \uncertain{tout} revêtement étale connexe $X_1$ de $X$, son image inverse $\overline{C}_1$ sur $\overline{C}$ est

\page{77}
connexe. \add{\uncertain{Il suffit de voir que} l'image inverse $C_1$ sur $C$ est géom.\ connexe ou \uncertain{même} géom.\ irréd.} Or $X_1$ provient d'un revêtement \struck{étale} \add{normal} intègre \struck{\ill{}} $X'_1 \to X'$ [EGA~V~\ill{}]. En vertu du th.\ de Bertini, l'image inverse \struck{\ill{}} dans $X'_1 \to \mathbb{P}^r$ \struck{\ill{}} de $L_t$ est \struck{géom.} géom.\ irréductible, ce qui implique que l'image $C_1$ \add{de \ill{}} sur $C$ \uncertain{l'est}, i.e.\ $\overline{C_1}$ est géom.\ connexe. Cela prouve le lemme.

3.3. De 3.2 résulte aussitôt que si \struck{l'extension} $C$ sur $K$ satisfait aux conclusions de 1.1, il en est de même de $X$ sur $k$. Compte tenu des réductions déjà faites dans 2.3, cela nous ramène, pour \uncertain{prouver} 1.1 \add{sur $k$}, au cas d'une courbe algébrique \add{\uncertain{lisse} sur $K$} déjà traité dans \struck{\ill{}} no~2. \uncertain{CQFD} \uncertain{Le} $k$ \uncertain{étant} de t.f., il en est de même de $K$\,!).

\textbf{Remarque 3.4.} Le lemme 3.2. Prouve aussi la \uncertain{même chose} que si $X$ \add{est} un schéma de t.f.\ sur un corps \add{\uncertain{de car.} $p \geq 0$} séparablement \uncertain{clos}, alors le \struck{\ill{}} plus grand \add{quotient} \add{\uncertain{premier} à $p$ $\pi_1(X,\xi)'$} \struck{$\mathbb{P}$} de $\pi_1(X,\xi)$, \uncertain{est} topologiquement \uncertain{de présentation} finie\add{\uncertain{\ill{}} (SGA~1~\ill{})}. En effet, la théorie de la \uncertain{descente} \struck{\ill{}} \uncertain{nous} permet de nous ramener au cas où $X$ est géom.\ unibranche \uncertain{normal} irréductible, \struck{et} puis au \struck{\ill{}} cas $X$ \struck{affine} \add{\ill{}}, puis 3.2. \struck{\ill{}} nous ramène au cas où $X$ est une \add{\ill{}} \uncertain{à une} courbe projective lisse $X'$. Or pour une telle $X$, \uncertain{le} \add{la} \uncertain{finitude} \uncertain{annoncée} du $\pi_1(X,\xi)'$ \uncertain{est} \uncertain{résultat} de (SGA~1~…).

\page{78}
\section*{4. Un complément au th.\ principal}

\textbf{Proposition 4.1.} Sous les conditions de 1.1, pour tout \struck{$k$} nb premier $\ell$ sauf un nb fini, \struck{le groupe} on a
\[
  \overline{[\pi', I(\ell)]} = I(\ell),
\]
i.e.\ le plus grand quotient \uncertain{premier} de $I(\ell)$ sur lequel $\pi_0$ opère trivialement est \emph{nul}.

On voit par dualité, en procédant comme dans 2.4, que cet énoncé équivaut au suivant :

\textbf{\struck{Proposition} Corollaire 4.2.} Pour presque tout $\ell$, on a
\[
  H^1(\overline{X}, \mathbb{Z}/\ell\mathbb{Z})^{\pi_0} = 0 .
\]

Pour prouver ceci, on est encore ramené, \struck{\ill{} \ill{}} par les réductions de 2.3, au cas où \struck{$X \subset X'$} $X$ est un ouvert d'un $X'$ projectif et géom.\ normal sur $k$. Dans ce cas, \struck{\ill{}} 3.2. nous ramène au cas où \uncertain{$X'$} est de dim.~1, donc est une courbe lisse et complète. \struck{\ill{} \ill{}} Prenant une extension finie de $k$, on peut supposer \uncertain{rationnels} \add{$x_i$} les pts de $X' - X$ \uncertain{rationnels} sur $k$. On utilise alors la suite exacte \uncertain{analogue} à (2.5.1)
\[
  \to H^1(\overline{X'}, \mathbb{Z}/\ell\mathbb{Z}) \to H^1(\overline{X}, \mathbb{Z}/\ell\mathbb{Z}) \to \prod_i H^2_{x_i}(\overline{X}, \mathbb{Z}/\ell\mathbb{Z})
\]
\[
  \to H^2(\overline{X'}, \mathbb{Z}/\ell\mathbb{Z}) \tag{4.2.1}
\]
qui nous montre que $H^1(\overline{X}, \mathbb{Z}/\ell\mathbb{Z})$ est une extension de $\left(\mu_\ell(\overline{k}^*)^{\otimes -1}\right)^{n-1}$ par $H^1(\overline{X'}, \mathbb{Z}/\ell\mathbb{Z})$. \struck{\ill{}} Soit
\[
  N = \text{card du groupe des racines de 1 dans } k^* \tag{4.2.2}
\]
\note{Le numéro de la première formule est surchargé ; on lit (4.2.1) ou (4.3.1), puis (4.2.2) ou (4.3.2).}

\page{79}
Alors, si \struck{\ill{}} $\ell \nmid N$, \struck{on aura $\mu_\ell(k^*)$} $\pi_0$ opère non trivialement sur le \add{$\mathbb{F}_\ell$-espace vectoriel} \struck{\ill{}} de dim~1 $\mu_\ell(\overline{k}^*)^{\otimes -1}$, donc sur son dual $V$, donc $V^{\pi_0} = 0$, ce qui implique que
\[
  H^1(\overline{X}, \mathbb{Z}/\ell\mathbb{Z})^{\pi} \subset H^1(\overline{X'}, \mathbb{Z}/\ell\mathbb{Z})^{\pi} \quad \text{si } \ell \nmid N . \tag{4.2.3}
\]
\struck{Or on a, si $A$ est la jacobienne de $X'$,}
\[
  \text{\struck{$H^1(\overline{X'}, \mathbb{Z}/\ell\mathbb{Z}) \simeq {}_\ell A(\overline{k})^\vee$}} ,
\]
\struck{où le $\vee$ désigne le dual : \ill{} les $\mathbb{Z}/\ell\mathbb{Z}$.}

D'autre part, \struck{$H^1(\overline{X'}, \mathbb{Z}/\ell\mathbb{Z})^\pi = 0$ signifie} comme dans 2.5 on se ramène, par la relation $H^1(\overline{X'}, \mathbb{Z}/\ell\mathbb{Z})^\pi = 0$ \add{pour presque tout $\ell$}, \struck{\ill{}} \struck{i.e.\ ${}_\ell A(\overline{k})^{\vee\pi} = 0$ i.e.\ $({}_\ell A(\overline{k}))_\pi = 0$} au cas où $k$ est fini. \struck{Alors \ill{} \ill{} Soit $f$ \ill{} \ill{} \ill{} bien connu, que} Soit $f$ le générateur \add{de Frobenius} de $\pi_0 \simeq \hat{\mathbb{Z}}$, et soit $V = {}_\ell A(\overline{k})$. Comme $V$ est fini, $\mathrm{Ker}(1-f_V)$ et $\mathrm{Coker}(1-f_V)$ ont \uncertain{même} cardinal, donc \struck{la relation}
\[
  ({}_\ell A(\overline{k}))_\pi = 0 \iff ({}_\ell A(\overline{k}))^\pi = 0 \quad \text{i.e.} \quad {}_\ell A(k) = 0 .
\]
Or soit $N'$ l'\struck{\ill{}} ordre du groupe \add{fini} $A(k)$. On voit que
\[
  {}_\ell A(k) = 0 \quad \text{si } \ell \nmid N' .
\]
Par suite, on trouve que si $\ell \nmid NN'$, on a
\[
  H^1(\overline{X'}, \mathbb{Z}/\ell\mathbb{Z})^\pi = 0, \quad \text{cqfd.}
\]
\note{Il écrit ici $\overline{X}^\ell$ ou $\overline{X'}$ ; la lecture $\overline{X'}$ suit l'argument.}

\textbf{Corollaire 4.3.} Soit $X$ un schéma \add{normal} géométriquement \struck{intègre} \add{\uncertain{loc.}\ de t.f.}\ sur un corps fini $\mathbb{F}_{p^r}$. Alors pour \uncertain{tout} nb premier $\ell \neq p$, l'homomorphisme $\pi_1(X,\xi)^{\mathrm{ab}}(\ell) \to \pi_1(\mathbb{F}_p) \simeq \hat{\mathbb{Z}}$ \struck{\ill{}} \add{a un noyau fini, et pour presque} tout $\ell$, c'est un isomorphisme.

\page{80}
\section*{5. Réduction du th.\ de monodromie pour $H^1(\overline{X}, \mathbb{Z}_\ell)$ au cas d'une courbe projective lisse géom.\ connexe}

5.1. La suite spectrale de \struck{\ill{}} descente pour un morphisme fini (SGA~4~VIII~…) nous ramène au cas $X$ \emph{normal}.

5.2. \struck{On peut alors supposer $X$ affine, et} Une extension \struck{\ill{}} \add{finie} du corps de base nous ramène au cas où de plus $X$ est un ouvert d'un schéma projectif \add{$X'$} \struck{\ill{}} géom.\ connexe et géom.\ \add{$t$} unibranche.

5.3. Dans ce cas, 3.2 (qui implique que
\[
  H^1(\overline{X}, \mathbb{Z}_\ell) \to H^1(\overline{C}, \mathbb{Z}_\ell)
\]
est injectif) nous ramène au cas \struck{où} $X'$ est \add{une} courbe \struck{\ill{}} projective et lisse.

[\emph{NB} \struck{Si $K = k(t_1, \ldots t_r)$, il est}, On choisit arbitrairement un anneau de valuation $V'$ du $K' = K(t)$ \add{\uncertain{\ill{}}} dominant $V$ ; dans l'image de $I' \to I$ est d'indice fini dans $V$]. Dans ce cas, \struck{la suite} on peut \add{encore} supposer que \add{les pts de} $X' - X$ sont rat.\ sur $K$. Alors, dans la suite exacte (2.5.1), l'opération de $\overline{I}$ sur les $H^2_{x}(\overline{X}, \mathbb{Z}_\ell) \simeq \mathbb{Z}_\ell(-1)$ est triviale, donc il \uncertain{revient} à prouver que \uncertain{l'action} sur $H^2(\overline{X'}, \mathbb{Z}_\ell)$ est quasi-unipotente, ce qui est la réduction annoncée.\note{Il écrit bien $H^2(\overline{X'},\mathbb{Z}_\ell)$ ; l'argument demande vraisemblablement $H^1$.}
\note{La section 5 se poursuit au-delà de ce lot.}

\end{document}
